\documentclass[11pt,a4paper]{article}
\usepackage[a4paper,margin=27mm,headheight=14pt]{geometry}
\usepackage{fontspec}
\setmainfont{Latin Modern Roman}
\setsansfont{Latin Modern Sans}
\setmonofont{Latin Modern Mono}
\usepackage{amsmath,amssymb,amsthm,mathtools,mathrsfs,enumitem,booktabs,array,longtable,fancyhdr}
\usepackage{xcolor}
\definecolor{MathInternalLink}{HTML}{1F4E79}
\definecolor{MathSectionLink}{HTML}{00656D}
\definecolor{MathEquationLink}{HTML}{643F7A}
\definecolor{MathCitationLink}{HTML}{2F6B3C}
\definecolor{MathExternalLink}{HTML}{843C39}
\usepackage{hyperref}
\hypersetup{colorlinks=true,linkcolor=MathInternalLink,citecolor=MathCitationLink,
urlcolor=MathExternalLink,linktoc=all,bookmarksnumbered=true,
pdftitle={Arithmetic Lefschetz theory, residual regulators, and integral heights},
pdfauthor={}}
\newtheorem{theorem}{Theorem}[section]
\newtheorem{proposition}[theorem]{Proposition}
\newtheorem{lemma}[theorem]{Lemma}
\newtheorem{corollary}[theorem]{Corollary}
\theoremstyle{definition}
\newtheorem{definition}[theorem]{Definition}
\newtheorem{conjecture}[theorem]{Conjecture}
\newtheorem{example}[theorem]{Example}
\newtheorem{question}[theorem]{Question}
\usepackage[nameinlink,noabbrev]{cleveref}
\crefname{conjecture}{conjecture}{conjectures}
\Crefname{conjecture}{Conjecture}{Conjectures}
\crefname{question}{question}{questions}
\Crefname{question}{Question}{Questions}
\DeclareRobustCommand{\MathSecref}[1]{\begingroup\hypersetup{linkcolor=MathSectionLink}\Cref{#1}\endgroup}
\DeclareRobustCommand{\MathEqref}[1]{\begingroup\hypersetup{linkcolor=MathEquationLink}\eqref{#1}\endgroup}
\newcommand{\Q}{\mathbb Q}\newcommand{\R}{\mathbb R}\newcommand{\C}{\mathbb C}
\newcommand{\Z}{\mathbb Z}\newcommand{\A}{\mathbb A}\newcommand{\Fbar}{\overline F}
\newcommand{\CH}{\operatorname{CH}}\newcommand{\Pic}{\operatorname{Pic}}
\newcommand{\Hom}{\operatorname{Hom}}\newcommand{\End}{\operatorname{End}}
\newcommand{\Spec}{\operatorname{Spec}}\newcommand{\Gal}{\operatorname{Gal}}
\newcommand{\rad}{\operatorname{rad}}\newcommand{\rk}{\operatorname{rank}}
\newcommand{\im}{\operatorname{im}}\newcommand{\Griff}{\operatorname{Griff}}
\newcommand{\AJ}{\operatorname{AJ}}\newcommand{\Tr}{\operatorname{Tr}}
\newcommand{\cl}{\operatorname{cl}}\newcommand{\ch}{\operatorname{ch}}
\newcommand{\Td}{\operatorname{Td}}\newcommand{\Supp}{\operatorname{Supp}}
\newcommand{\homol}{\mathrm{hom}}\newcommand{\alg}{\mathrm{alg}}
\newcommand{\et}{\textnormal{ét}}\newcommand{\cont}{\mathrm{cont}}
\newcommand{\cP}{\mathcal P}\newcommand{\cO}{\mathcal O}
\newcommand{\cD}{\mathscr D}\newcommand{\cV}{\mathcal V}
\newcommand{\tGamma}{{}^t\Gamma}
\numberwithin{equation}{section}
\setlist{leftmargin=2em,itemsep=3pt,topsep=5pt}
\setlength{\emergencystretch}{2em}
\setlength{\parskip}{3pt}
\pagestyle{fancy}\fancyhf{}
\fancyhead[L]{\small Lefschetz theory and integral heights}
\fancyhead[R]{\small Research archive}\fancyfoot[C]{\thepage}
\title{Arithmetic Lefschetz theory, residual regulators,\\
and integral heights on products of curves}
\author{}
\date{Research note, second round --- 16 November 2026}
\begin{document}
\maketitle
\begin{abstract}
For a product of pointed curves, we identify noncentral vanishing in
the all-\(h^1\) Chow groups of its subproducts with arithmetic hard
Lefschetz for product polarizations. An explicit primitive identity
expresses the signed height as a weighted sum of central heights.
These identities reduce primitive integral Northcott to the central
algebraic Abel--Jacobi kernel and residual integral Northcott.
We construct a residual arithmetic Abel--Jacobi map on the central
Griffiths quotient by applying an actual Chow projector to continuous
étale cohomology. We then separate integral Northcott finiteness into
three arithmetic problems: rational detection, a common integral
lattice, and control of the fields of definition. Trace and
Hermite--Minkowski give an exact ramification criterion for the last
problem. Finally, in Jacobian twisting families we prove a normalized
Quillen projection-formula isometry and its rank-zero Deligne-pairing
specialization. The height statements apply to a given system with
correspondence adjunction and the classical curve normalization.
\end{abstract}
{\hypersetup{linkcolor=MathSectionLink}\tableofcontents}

\section{Introduction}\label{sec:introduction}

Let \(F\) be a number field. For a smooth projective variety \(X/F\)
of dimension \(n\), set
\[
 V^r(X_E)=\CH^r(X_E)_{\homol,\Q}
\]
for every finite extension \(E/F\). A Beilinson--Bloch height pairs
\(V^r\) with \(V^{n+1-r}\). Nondegeneracy, positivity and finiteness
of integral height balls ask different questions about this pairing.
On products of curves, the tensor Chow projectors make their
relationship explicit.

Write \(X=\prod_{i=1}^n C_i\), choose points \(o_i\in C_i(F)\), and
put \(H_i=p_i^*[o_i]\). For \(S\subseteq\{1,\ldots,n\}\), let
\[
 Y_S=\prod_{i\in S}C_i,\qquad
 \Pi_S=\bigotimes_{i\in S}
  (\Delta_{C_i}-o_i\times C_i-C_i\times o_i),\qquad
 U_S^c(E)=\Pi_{S*}V^c(Y_{S,E}).
\]
The empty subproduct has zero homologically trivial Chow group.
The central degree of a nonempty subproduct is characterized by
\(|S|=2c-1\).

\begin{theorem}[Arithmetic Lefschetz reduction]
\label[theorem]{thm:main-lefschetz}
Fix positive integers \(a_i\) and
\(\ell=\sum_i a_iH_i\). For every finite \(E/F\), the following
conditions are equivalent:
\begin{enumerate}[label=(\roman*)]
\item \(U_S^c(E)=0\) for every \(S,c\) with \(|S|\ne2c-1\).
\item For \(1\le r\le\lfloor(n+1)/2\rfloor\),
\[
 \ell^{n+1-2r}:V^r(X_E)\xrightarrow{\sim}V^{n+1-r}(X_E).
\]
\end{enumerate}
In particular, condition (ii) for one choice of positive integers
\(a_i\) is equivalent to condition (ii) for every such choice,
with the points \(o_i\) fixed.
\end{theorem}

The proof in \MathSecref{sec:lefschetz-proof} is a calculation in
actual Chow groups. Each summand is a tensor product
\[
 U_S^c(E)\otimes_\Q
 \Q[x_i:i\notin S]/(x_i^2),
\]
where \(x_i\) represents \(a_iH_i\). Multiplication by \(\ell\)
acts on the second factor. A noncentral core places the centre of
this Lefschetz string in the wrong degree, forcing a failure of
injectivity or surjectivity.

The height-theoretic statements use a given system of pairings
\(h_X\) with the four properties in
\MathSecref{sec:height-data}: symmetry and bilinearity, absolute
field-extension invariance, correspondence adjunction, and the
negative Néron--Tate normalization on curves. Under these assumptions,
the same calculation yields an exact identity. If \(|S|=2c-1\),
\(m=n-|S|\), \(t=r-c\), and
\(z_S=\sum_{|T|=t}u_Tx_T\) is primitive, then
\begin{equation}\label{eq:intro-squares}
 (-1)^r h_X(z_S,\ell^{n+1-2r}z_S)
  =(m-2t)!\prod_{i\notin S}a_i
       \sum_{|T|=t}(-1)^c h_{Y_S}(u_T,u_T).
\end{equation}
Here \(x_T=\prod_{i\in T}x_i\), interpreted by the fixed Chow
embeddings of \MathSecref{sec:chow-blocks}. The proof in
\MathSecref{sec:primitive-proof} determines both the sign and the
factorial. It reduces real primitive positivity and integral
Northcott on \(X\) to central cores, with explicit constants and
uniform denominators.

The central core contains two different arithmetic pieces. Suppose
\(Y=\prod_{i=1}^{2c-1}C_i\), \(U_E=\Pi_*V^c(Y_E)\), and
\(A_E=U_E\cap\CH^c(Y_E)_{\alg,\Q}\), where algebraic triviality is
geometric. After a fixed finite extension \(F_0/F\), an actual
self-transpose projector \(e\le\Pi\) realizes the maximal
curve-generated subspace of \(\Pi H^{2c-1}(Y)\). Put \(R=\Pi-e\).
The construction and its consequences are proved in
\MathSecref{sec:central-projector}. They give
\[
 U_E=eU_E\oplus RU_E,\qquad
 RA_E=\ker(\AJ_\Q|_{A_E}),\qquad
 RU_E/RA_E\simeq G_E:=U_E/A_E.
\]
The first summand has a classical Néron--Tate height. Orthogonalizing
by \(R\) gives a canonical residual pairing \(h_G\) on \(G_E\).

For the projected integral images, central Northcott is equivalent
to vanishing of \(RA_E\) and residual Northcott
(\Cref{thm:central-northcott}). Northcott itself forces real
positive definiteness, even without a finite-rank assumption:
rounding multiples of a real null vector would otherwise give
infinitely many integral vectors of bounded height
(\Cref{lem:northcott-real}). Together with the primitive formula,
these facts give the exact equivalence between the product and
residual programs in \Cref{thm:program-equivalence}.

\begin{theorem}[Residual arithmetic regulator]
\label[theorem]{thm:main-regulator}
For every prime \(\ell_0\) and finite \(E/F_0\), set
\[
 V_{\ell_0}=R_*H^{2c-1}_{\et}(Y_{\Fbar},\Q_{\ell_0}(c)).
\]
There is a natural \(\Q\)-linear map
\[
 a_{\ell_0,E}:G_E\longrightarrow H^1(E,V_{\ell_0})
\]
induced by the continuous étale cycle class. It is independent of
the curve-visible Chow lift under the common cohomological
identification, and is compatible with restriction and trace.
There is one finite set \(S_0\) of places of \(F_0\) such that
inflation identifies
\[
 H^1(G_{E,S_E},V_{\ell_0})\simeq H^1(E,V_{\ell_0}),
\]
where \(S_E\) consists of places above \(S_0\) and \(\ell_0\).
This whole target is finite-dimensional over \(\Q_{\ell_0}\).
\end{theorem}

In \MathSecref{sec:regulator-proof}, the \(R\)-part of
Hochschild--Serre has exactly one geometric row. A curve-generated
algebraically trivial class maps to zero on its only associated
graded term, hence to zero in the whole target. This provides an
arithmetic detector on the Griffiths quotient without choosing a
splitting of a multi-step filtration. The proof is completed in
\MathSecref{sec:fixed-model}: the residual representation has
weight \(-1\), so local Galois cohomology at good places proves
finite-dimensionality of the whole target. Scalar-extended
injectivity of the regulator would imply finite rational rank.

Finite rank and real positivity still leave an integral problem.
For a subgroup \(\Lambda\) of a finite-dimensional rational space
with a real positive quadratic form, Northcott is equivalent to one
uniform denominator bound in a rational basis. We prove an
all-prime saturation criterion for this bound in
\MathSecref{sec:all-prime}. We also give a weak-Mordell--Weil plus
height-gap criterion, including an explicit bound for the number
of classes in a height ball.

The field of definition introduces a further condition. If an
integral class \(\xi\) is represented over \(E\), its minimal
\emph{class} field \(K=F_0(\xi)\) may be smaller than the field of
its representative. Trace nevertheless gives
\[
 [E:K]\xi\in\Lambda(K).
\]
Consequently, fixed-field Northcott together with a finite
ramification set for minimal class fields at bounded height and
degree is equivalent to bounded-degree Northcott
(\MathSecref{sec:ramification-proof}). An explicit Galois-module
example shows why uniform bounds on ranks, height gaps and Kummer
quotients alone do not settle the varying-field problem.

Finally, \MathSecref{sec:metrics} treats a metric comparison on a
proper parameter space. Let \(C\) be a pointed complex curve,
\(J=\Pic^0(C)\), and \(\cP\) the normalized Poincaré bundle on
\(C\times J\). Twisting a fixed perfect complex on \(C\times X\)
by \(\cP\) gives a normalized Quillen determinant line on \(J\).
Derived projection to \(C\) preserves this normalized metric
exactly. In virtual rank zero, it is the flat Deligne-pairing
metric. The proof combines the Bismut--Gillet--Soulé curvature
formula with the origin rigidification on the compact Jacobian.

These reductions fit within established conjectural frameworks.
The arithmetic Lefschetz and primitive positivity assertions are
the Beilinson conjectures recorded in
\cite[Conjecture~2.3.1]{ZhangStandard}. The arithmetic regulators use
the continuous-cohomology framework of
\cite{Raskind,Kahn}; the determinant argument uses
\cite{BGS,BiswasSchumacher}. The final section formulates the
remaining detection, saturation and ramification problems for the
specific residual motive. The proofs below supply reductions and
comparison statements within these frameworks.

\section{Height data and the central Chow projector}\label{sec:setup}

\subsection{Height data and integral images}\label{sec:height-data}

Fix an algebraic closure \(\Fbar\) of \(F\) and an embedding into
\(\C\). Homological triviality is defined by geometric rational
Betti cohomology. All Chow groups below have rational coefficients
unless otherwise indicated. We regard \(V^r(X_E)\) as a subspace
of \(V^r(X_{\Fbar})\): restriction is injective because a rational
equivalence descends to a finite extension and trace composed with
restriction multiplies by its degree.

\begin{definition}\label[definition]{def:heights}
A curve-normalized height system consists of pairings
\[
 h_X:V^r(X_E)\times V^{n+1-r}(X_E)\longrightarrow\R
\]
over every finite extension \(E/F\), satisfying:
\begin{enumerate}[label=(H\arabic*)]
\item bilinearity over \(\Q\) and symmetry under exchange of the
complementary degrees;
\item invariance under extension in the absolute normalization,
and covariance under isomorphisms, including Galois conjugation;
\item correspondence adjunction
\[
 h_X(\Gamma_*z,w)=h_Y(z,{}^t\Gamma_*w);
\]
\item on a pointed curve \(C\), the form \(-h_C\) on
\(\Pic^0(C)_\Q\) is the positive Néron--Tate form for the canonical
principal polarization, in the corresponding Poincaré normalization.
\end{enumerate}
\end{definition}

These hypotheses will be understood in every statement involving
\(h\). A real positive form means positive definiteness after
tensoring the rational group over a fixed number field with \(\R\).
For codimension \(r\) on \(X\), the symbol \(\AJ_\Q\)
denotes the complex Abel--Jacobi map with target
\(J^r(X_\C)(\C)\otimes_\Z\Q\).

Let \(\Lambda_X^r(E)\) be the image in \(V^r(X_{\Fbar})\) of the
integral-coefficient cycles over \(E\) with zero geometric rational
cycle class. Thus torsion and rational equivalence are removed before
counting. For a fixed rational projector \(P\), we use its projected
integral image \(P_*\Lambda_X^r(E)\); for an integral projector this
equals \(\Lambda_X^r(E)\cap P_*V^r(X_E)\).

\begin{definition}\label[definition]{def:northcott}
For compatible integral images \(\Lambda(E)\) in one geometric
rational group and a quadratic form \(q\), bounded-degree Northcott
means that, for every \(B\in\R\) and integer \(D\ge1\),
\[
 \{\xi\in\bigcup_{[E:F]\le D}\Lambda(E):q(\xi)\le B\}
\]
is finite. Fixed-field Northcott means the same assertion with
\(E\) fixed. Bounded denominators mean the union of
\(d^{-1}\Lambda(E)\) for \(1\le d\le d_0\), with \(d_0\) fixed.
\end{definition}

Northcott for integral images implies the bounded-denominator version
by scaling. A fixed initial extension \(F_0/F\) is harmless for
bounded-degree assertions: after replacing \(E\) by \(EF_0\),
\([EF_0:F_0]\le[E:F]\), the integral image remains integral, and
the absolute height is unchanged.

\subsection{Curve-invisible cycles}\label{sec:invisible-cycles}

\begin{lemma}\label[lemma]{lem:invisible}
Let \(C\) be a smooth projective curve and
\(\Gamma\in\CH^r(C\times X)_\Q\), with \(\dim X=n\). If
\(\Gamma_*H^1(C)=0\) in \(H^{2r-1}(X)\), then
\[
 {}^t\Gamma_*w=0\quad\text{in }\Pic^0(C)_\Q
 \qquad(w\in V^{n+1-r}(X)).
\]
Hence \(\Gamma_*\Pic^0(C)_\Q\) is in the left height radical.
\end{lemma}

\begin{proof}
By Poincaré duality, the transpose map from
\(H^{2n+1-2r}(X)\) to \(H^1(C)\) is zero. Functoriality of the
Abel--Jacobi map under Chow correspondences therefore makes the
Abel--Jacobi class of \({}^t\Gamma_*w\) zero. On a curve, the
rationalized Abel--Jacobi map is the identification with
\(\Pic^0(C)_\Q\). Vanishing after complex base change descends for
divisor classes with rational coefficients. The final assertion is
adjunction. The correspondence functoriality used here is also
proved in \cite[Lemma~4.11]{LagaShnidman}.
\end{proof}

Algebraically trivial cycles are generated by images of degree-zero
divisors on smooth projective curves. For geometric algebraic
triviality over \(E\), the finitely many data can be defined over a
finite extension and then traced to \(E\). This convention is used
throughout.

\subsection{The maximal curve-visible projector}
\label{sec:central-projector}

Let \(k=2c-1\), \(Y=\prod_{i=1}^kC_i\), and
\(\Pi=\bigotimes_i\pi_{1,i}\). Set
\[
 U_E=\Pi_*V^c(Y_E),\qquad
 A_E=U_E\cap\CH^c(Y_E)_{\alg,\Q},\qquad G_E=U_E/A_E.
\]
For the product polarization \(\ell_Y=\sum a_i[o_i]\), the actual
Chow identity \(\ell_Y\Pi=0\) follows from
\([o_i]\pi_{1,i}=0\). Thus the geometric cohomology of \(\Pi\)
is middle and primitive.

Let \(W\subset\Pi H^k(Y,\Q)\) be the span of all images
\[
 H^1(C)\longrightarrow\Pi H^k(Y,\Q)
\]
of codimension-\(c\) curve correspondences over finite extensions
of \(F\). Choose finitely many correspondences spanning \(W\),
and one field \(F_0\) defining them and their source points.
Write their matrix as \(\Gamma=(\Gamma_j)_j\), with
\(\Gamma_j=\Pi\Gamma_j\pi_{1,j}\), and put
\(N=\bigoplus_jh^1(C_j)\), \(J=\prod_j\Pic^0(C_j)\).

\begin{proposition}\label[proposition]{prop:central-projector}
There are actual Chow projectors \(e,R\) over \(F_0\), with
\[
 e={}^te,\quad e^2=e,\quad e\Pi=\Pi e=e,\quad R=\Pi-e,
\]
such that the following hold for every finite \(E/F_0\).
\begin{enumerate}[label=(\roman*)]
\item The cohomological action of \(e\) is the orthogonal
projection onto \(W\). There is a fixed abelian variety \(B/F_0\)
with \(eU_E\simeq B(E)_\Q\), through fixed Chow correspondences.
The form \(q_e=(-1)^ch_Y\) on this summand is a positive
Néron--Tate form.
\item With \(K_E=RA_E\),
\[
 K_E=\ker(\AJ_\Q|_{A_E})
    =A_E\cap\rad(h_Y|_{U_E})
    =\rad(h_Y|_{A_E}).
\]
\item There is an orthogonal decomposition
\(U_E=eU_E\oplus RU_E\). The residual form on
\(RU_E/RA_E\simeq G_E\) is independent of the choice of maximal
curve-visible projector. Write it as \(h_G\) and set
\(q_G=(-1)^ch_G\).
\end{enumerate}
\end{proposition}

\begin{proof}
Set \(a={}^t\Gamma\Gamma\). Since \(2c-k=1\), its entries are
divisor correspondences, so \(a\in\End(N)=\End(J)_\Q\).
This realization is faithful, and transposition corresponds to
Rosati adjunction for the product principal polarization.
Geometric Hodge--Riemann on the primitive subspace \(W\) implies
that \((-1)^{c-1}a\) is positive semidefinite and
\(\ker a=\ker\Gamma\) on \(H^1\). In particular \(a\) is
semisimple. Choose a polynomial \(f\in\Q[t]\) equal to
\(t^{-1}\) on the nonzero spectrum and zero at zero, and set
\[
 a^+=f(a),\qquad p=a^+a,\qquad
 e=\Gamma a^+{}^t\Gamma.
\]
Faithfulness on \(N\) makes \(a^+aa^+=a^+\), \(p^2=p={}^tp\)
actual Chow identities. These identities prove the assertions about
\(e\) and its cohomology. Let \(B=pJ\), up to isogeny.
The maps \(\Gamma\) and \(a^+{}^t\Gamma\) identify
\(B(E)_\Q\) with \(eU_E\). Adjunction gives
\[
 h_Y(\Gamma x,\Gamma y)=h_{\mathrm{curves}}(x,ay).
\]
The curve normalization and the positivity of
\((-1)^{c-1}a\) on \(pJ\) give (i).

For every later curve correspondence \(\Delta\), the geometric
\(H^1\) action of \(R\Delta\) is zero by maximality of \(W\).
Curve generation and \Cref{lem:invisible} give
\(RA_E\subseteq\rad(h_Y|_{U_E})\). The Abel--Jacobi map is zero
on \(RA_E\) by the same geometric action. Its restriction to
\(eU_E\) is injective after rationalization: the corresponding
map from \(B\) to the intermediate Jacobian has injective rational
\(H_1\), hence finite kernel. Since \(eU_E\subseteq A_E\) and
its height is nondegenerate, all equalities in (ii) follow.

The self-transpose projectors give the orthogonal decomposition.
The radical assertion just proved makes the residual form well
defined on \(G_E\). Two orthogonalized lifts of the same Griffiths
class differ by an element of \(A_E\) orthogonal to \(A_E\).
By (ii), their difference is in the full radical. Thus the residual
form is independent of the lift and of \(e\).
\end{proof}

The cohomological subspace \(W\) and its orthogonal complement are
defined independently of the chosen spanning correspondences. Betti
and étale comparison transports them to each \(\Q_{\ell_0}\)
realization. The algebraic construction of \(e\) in the proof uses
geometric Hodge--Riemann and the faithful realization of \(h^1\)
endomorphisms; it is available independently of the height data.

\begin{lemma}\label[lemma]{lem:visible-integral}
The projected integral image in \(eU_E\) has bounded-degree
Northcott for \(q_e\), allowing any fixed denominator bound.
\end{lemma}

\begin{proof}
Put \(D=a^+{}^t\Gamma\). These fixed curve correspondences satisfy
\(De=D\). One integer \(N\) clears their denominators, so an
integral cycle \(z\) gives a genuine point \(NDz\in J(E)\).
The same point detects \(ez\), even though \(ez\) may have
rational coefficients.

On the fixed factor \(pJ(\Fbar)_\Q\), the ambient Néron--Tate
height of \(J\) is bounded above by a fixed positive multiple of
\(q_e\) under the inverse \(D\). This is comparison of two
polarizations on \(B\); passing to rational points removes
torsion without changing either height. Thus a bound for
\(q_e(ez)\) bounds the height of the actual point \(NDz\).
Classical bounded-degree Northcott on \(J\), followed by the
injectivity of \(D\) on \(eU_E\), proves the assertion. Fixed
extra denominators are absorbed into \(N\). See
\cite[Sections~9--10]{ConradFeng} for the height and Mordell--Weil
facts used here.
\end{proof}

\section{Lefschetz strings in the Chow group}\label{sec:lefschetz}

\subsection{Exact Chow blocks}\label{sec:chow-blocks}

On a pointed curve put
\[
 \pi_0=o\times C,\qquad \pi_2=C\times o,\qquad
 \pi_1=\Delta_C-\pi_0-\pi_2.
\]
They are orthogonal projectors, with
\({}^t\pi_0=\pi_2\) and \({}^t\pi_1=\pi_1\). For a partition
\([n]=S\sqcup T\sqcup U\), define
\[
 P_{S,T,U}=\bigotimes_{i\in S}\pi_{1,i}
            \otimes\bigotimes_{i\in T}\pi_{2,i}
            \otimes\bigotimes_{i\in U}\pi_{0,i}.
\]
Let \(j_T\) insert the point \([o_i]\) on factors in \(T\) and
the fundamental class \([C_i]\) on factors in \(U\). Its inverse
on this projector image is the fixed correspondence
\[
 b_T(z)=\Pi_{S*}p_{S*}
           \left(z\cdot\prod_{i\in U}H_i\right).
\]
Indeed \(b_Tj_T=\Pi_S\) and \(j_Tb_T=P_{S,T,U}\). Therefore
\begin{equation}\label{eq:chow-decomposition}
 V^r(X_E)=\bigoplus_{S,T,U}j_TU_S^{r-|T|}(E).
\end{equation}
These identities are also valid on integral images. All intersections
and pushforwards in \(b_T\) have integral coefficients.

Put \(I=[n]\setminus S\), \(m=|I|\), and
\[
 B_I=\Q[x_i:i\in I]/(x_i^2),\qquad \deg x_i=1.
\]
For \(T\subset I\), write \(a_T=\prod_{i\in T}a_i\). The
identification
\begin{equation}\label{eq:weighted-chow-block}
 u\otimes x_T\longmapsto a_Tj_Tu
\end{equation}
groups \MathEqref{eq:chow-decomposition} as
\[
 V^*(X_E)=\bigoplus_{S,c} U_S^c(E)\otimes B_I[c],
\]
where a vector in \(U_S^c\otimes B_I^t\) has codimension \(c+t\).
Multiplication by \(\ell=\sum_i a_iH_i\) becomes multiplication
by \(E_I=\sum_{i\in I}x_i\) on the second factor. On the other
factors it vanishes because \(H_i\pi_{1,i}=0\) as a Chow
correspondence. This last identity follows by composing
\(\mu_{[o_i]}=[o_i]\times[o_i]\) with the three curve projectors.

If \(z\) is integral, its coefficient \(u_T\) in
\MathEqref{eq:weighted-chow-block} satisfies
\(a_Tu_T\in\Lambda_{U_S^c}(E)\). Thus the single integer
\begin{equation}\label{eq:coordinate-denominator}
 N_a=\prod_{i=1}^n a_i
\end{equation}
clears every coefficient denominator, in every degree and over every
finite extension.

\subsection{A square-zero Lefschetz calculation}\label{sec:square-zero}

On \(B_I\), use the monomial basis \(x_T\) and define
\[
 E x_T=\sum_{i\notin T}x_{T\cup\{i\}},\qquad
 F x_T=\sum_{i\in T}x_{T\setminus\{i\}},\qquad
 Hx_T=(m-2|T|)x_T.
\]
Then \([F,E]=H\), \([H,F]=2F\), and \([H,E]=-2E\).
The operators \(E,F\) are adjoint for the positive inner product
in which the monomials are orthonormal. Let
\(*x_T=x_{I\setminus T}\).

\begin{lemma}\label[lemma]{lem:square-zero}
For \(0\le t\le m/2\),
\[
 E^{m-2t}:B_I^t\xrightarrow{\sim}B_I^{m-t},\qquad
 \ker(E^{m-2t+1}|_{B_I^t})=\ker(F|_{B_I^t}).
\]
For every rational vector space \(M\) and every
\(z\in M\otimes\ker(F|_{B_I^t})\),
\begin{equation}\label{eq:star-identity}
 E^{m-2t}z=(-1)^t(m-2t)!\,*z.
\end{equation}
\end{lemma}

\begin{proof}
The displayed commutators give the standard finite-dimensional
\(\mathfrak{sl}_2\) representation on \(B_I\). More explicitly,
if \(Fv=0\), \(v\in B_I^t\), then
\[
 FE^jv=j(m-2t-j+1)E^{j-1}v.
\]
Adjointness shows that the string generated by a nonzero such \(v\)
has nonzero terms through \(E^{m-2t}v\) and then terminates.
Decomposing successively into \(\ker F\) and \(\im E\), which
are orthogonal, gives the Lefschetz strings and the first two
assertions.

For the exact constant, write \(z=\sum_{|T|=t}u_Tx_T\).
From \(Fz=0\) one obtains, for every \(J\subset I\) with
\(|J|<t\),
\[
 \sum_{\substack{T\supset J\\|T|=t}}u_T=0.
\]
This follows by looking at the coefficient of \(x_J\) in
\(F^{t-|J|}z\), which is \((t-|J|)!\) times the sum. If
\(K\subset I\) has size \(m-t\) and \(R_0=I\setminus K\),
inclusion--exclusion gives
\[
 \sum_{\substack{T\subset K\\|T|=t}}u_T
 =\sum_{J\subset R_0}(-1)^{|J|}
       \sum_{\substack{T\supset J\\|T|=t}}u_T
 =(-1)^tu_{R_0}.
\]
The coefficient of \(x_K\) in \(E^{m-2t}z\) is
\((m-2t)!\) times the left side. This proves
\MathEqref{eq:star-identity}, including \(t=0\).
\end{proof}

\subsection{Proof of the arithmetic Lefschetz reduction}
\label{sec:lefschetz-proof}

We prove \Cref{thm:main-lefschetz}. Suppose first that all the
noncentral cores vanish. On a remaining block put \(k=|S|=2c-1\),
\(m=n-k\), and \(t=r-c\). Then
\[
 n+1-2r=m-2t,
\]
and the desired map is
\(1\otimes E^{m-2t}\). It is an isomorphism by
\Cref{lem:square-zero}. If \(t<0\), both its source and target
in this block vanish. Taking the direct sum proves (ii).

Conversely, suppose \(U_S^c(E)\ne0\) and put
\(d=k+1-2c\). The homologically trivial groups force
\(1\le c\le k\). If \(d>0\), take \(r=c\). The required
map has exponent \(m+d>m\), so it kills
\(U_S^c(E)\otimes1\); it cannot be injective. If \(d<0\),
take \(r=k+1-c\). This satisfies
\(1\le r\le\lfloor(n+1)/2\rfloor\). The target has the
nonzero component \(U_S^c(E)\otimes x_I\) in degree \(c+m\),
whereas the source in this same block has polynomial degree
\(r-c=d<0\) and is zero. Multiplication by \(\ell\) preserves
the block, so surjectivity fails. Hence \(d=0\) for every nonzero
core. This proves (i).

\subsection{Primitive heights and their exact constants}
\label{sec:primitive-proof}

Adjunction and \({}^tP_{S,T,U}=P_{S,U,T}\) show that Chow blocks
are height-orthogonal unless their point and fundamental-class
indices are complementary and their \(S\) indices agree. In the
complementary case,
\begin{equation}\label{eq:complement-height}
 h_X(j_Tu,j_{I\setminus T}v)=h_{Y_S}(u,v).
\end{equation}
Indeed the transpose of the first embedding composed with the
second integrates one degree-one point in each omitted factor.

Define
\[
 P_\ell^r(X_E)=\ker\bigl(
   \ell^{n+2-2r}:V^r(X_E)\to V^{n+2-r}(X_E)\bigr)
\]
for \(1\le r\le\lfloor(n+1)/2\rfloor\).

\begin{theorem}[Primitive sum of squares]
\label[theorem]{thm:primitive-squares}
Let \(|S|=2c-1\), \(m=n-|S|\), and \(0\le t=r-c\le m/2\).
Write a vector in the corresponding block as
\[
 z_S=\sum_{|T|=t}u_Tx_T,
 \qquad u_T\in U_S^c(E).
\]
It is primitive exactly when \(Fz_S=0\). For such a vector,
\begin{equation}\label{eq:primitive-squares}
 q_{\ell,r}(z_S):=(-1)^rh_X(z_S,\ell^{n+1-2r}z_S)
   =(m-2t)!\,a_I\sum_{|T|=t}q_{S,c}(u_T),
\end{equation}
where \(a_I=\prod_{i\in I}a_i\) and
\(q_{S,c}(u)=(-1)^ch_{Y_S}(u,u)\).
\end{theorem}

\begin{proof}
The primitive exponent is \(m-2t+1\), so the equivalence follows
from \Cref{lem:square-zero}. The same lemma gives
\(\ell^{n+1-2r}z_S=(-1)^t(m-2t)!\,*z_S\).
In pairing \(z_S\) with \(*z_S\), only the identical coefficient
indices survive by \MathEqref{eq:complement-height}. Their weights
multiply to \(a_Ta_{I\setminus T}=a_I\). The remaining sign is
\((-1)^{r+t}=(-1)^c\), since \(r=c+t\). This proves the formula.
\end{proof}

\begin{corollary}\label[corollary]{cor:positivity-propagation}
Suppose all the noncentral subproduct cores vanish and all central
forms \(q_{S,c}\) are real positive definite over a fixed number
field. Then the complementary height pairings on \(X\) are
nondegenerate over that field, also after real extension, and
\(q_{\ell,r}\) is real positive definite on \(P_\ell^r\).
Its associated bilinear form \(b_{\ell,r}\) satisfies
\[
 |b_{\ell,r}(z,w)|^2
      \le q_{\ell,r}(z)q_{\ell,r}(w).
\]
If noncentral vanishing holds over every finite extension and each
central integral image has bounded-degree Northcott, then so does
\(\Lambda_X^r(E)\cap P_\ell^r\) for \(q_{\ell,r}\).
\end{corollary}

\begin{proof}
The decomposition into complementary blocks and the definiteness
of their central forms prove nondegeneracy. Distinct \(S\) are
orthogonal. Formula \MathEqref{eq:primitive-squares} therefore
proves primitive definiteness. Cauchy--Schwarz follows by applying
positivity to \(q_{\ell,r}(z+tw)\), \(t\in\R\).

For the bounded-degree assertion, central Northcott gives
nonnegativity over each finite extension: a nonzero rational class
of nonpositive height would have infinitely many integral multiples
in a bounded height ball. For an integral primitive \(z\) of
height at most \(B\), every
coefficient satisfies
\[
 q_{S,c}(u_T)\le
   \frac{B}{(m-2t)!a_I},\qquad
 N_a u_T\in\Lambda_{U_S^c}(E).
\]
Thus central Northcott and \MathEqref{eq:coordinate-denominator}
give finitely many possibilities for every coefficient. There are
only finitely many blocks, and the decomposition is injective.
\end{proof}

\section{The exact central Northcott reduction}
\label{sec:central-northcott}

Continue with the central notation of
\MathSecref{sec:central-projector}. Let
\(\Lambda_U(E)=\Pi_*\Lambda_Y^c(E)\) and let
\(\Lambda_G(E)\) be its image in \(G_E\). The quotient may be
viewed in the geometric group: if an element defined over \(E\)
becomes geometrically algebraically trivial, it is in \(A_E\) by
definition. In particular, these rational quotients inject under
field extension.

\begin{theorem}\label[theorem]{thm:central-northcott}
The following conditions are equivalent:
\begin{enumerate}[label=(\roman*)]
\item The form \(q_U(z)=(-1)^ch_Y(z,z)\) has bounded-degree
Northcott on \(\Lambda_U\).
\item For every finite \(E/F_0\), \(K_E=\ker(\AJ_\Q|_{A_E})=0\),
and \(q_G\) has bounded-degree Northcott on \(\Lambda_G\).
\end{enumerate}
There is also a fixed-field version, with all groups and the
vanishing condition taken over one fixed \(E\).
\end{theorem}

\begin{proof}
Suppose (i). If \(0\ne z\in K_E\), multiply by one integer to
make \(z\) integral. All positive integer multiples of this class
are distinct and have height zero, contradicting Northcott.
Thus \(K_E=0\).

Choose an integer \(M\ge1\) such that \(MR\) is represented by
an integral correspondence. For \(\xi\in\Lambda_G(E)\), choose
\(z\in\Lambda_U(E)\) mapping to \(\xi\). Then
\[
 MRz\in\Lambda_U(E),\qquad
 q_U(MRz)=M^2q_G(\xi),\qquad [MRz]=M\xi.
\]
Northcott for \(\Lambda_U\) bounds the possible classes \(MRz\),
and hence the possible \(\xi\). This proves residual Northcott.

Conversely, assume (ii). Residual Northcott implies
\(q_G(\xi)>0\) for every nonzero rational class \(\xi\in G_E\).
Indeed one multiple of \(\xi\) is in \(\Lambda_G(E)\), and a
nonpositive height would give infinitely many integral multiples
of height at most zero. For every \(z\in U_E\), orthogonality gives
\begin{equation}\label{eq:central-norm-split}
 q_U(z)=q_e(ez)+q_G([z]).
\end{equation}
Both terms are nonnegative. If \(z\) is integral and
\(q_U(z)\le B\), residual Northcott bounds the possibilities for
\([z]\). The denominators of the fixed projector \(e\), followed
by \Cref{lem:visible-integral}, bound the possibilities for \(ez\).
Finally
\[
 z\longmapsto(ez,[z])
\]
is injective: its kernel is \(RA_E=K_E=0\). This proves (i).
The argument with \(E\) fixed is identical.
\end{proof}

The equivalence separates the algebraic Abel--Jacobi kernel from
integral finiteness on the residual quotient. In fact, Northcott
also forces real positive definiteness, by
\Cref{lem:northcott-real}. This observation will combine the
central equivalence with the primitive calculation in
\Cref{thm:program-equivalence}.

\section{Residual arithmetic Abel--Jacobi maps}\label{sec:regulators}

\subsection{The single-row construction}\label{sec:regulator-proof}

We use continuous étale cohomology and its cycle class map as in
\cite{Raskind,Kahn}. Fix a rational prime \(\ell_0\), and put
\[
 k=2c-1,\qquad
 V_{\ell_0}=R_*H^k_{\et}(Y_{\Fbar},\Q_{\ell_0}(c)).
\]
The action of the actual Chow idempotent \(R\) on geometric
cohomology is concentrated in degree \(k\). Apply it to the
correspondence-compatible Hochschild--Serre spectral sequence
\[
 H^p\bigl(E,H^q_{\et}(Y_{\Fbar},\Q_{\ell_0}(c))\bigr)
 \Longrightarrow H^{p+q}_{\cont}(Y_E,\Q_{\ell_0}(c)).
\]
Only the row \(q=k\) survives. Hence there is a canonical
identification
\begin{equation}\label{eq:single-row}
 R_*H^{2c}_{\cont}(Y_E,\Q_{\ell_0}(c))
       \simeq H^1(E,V_{\ell_0}).
\end{equation}
On this summand the Hochschild--Serre filtration satisfies
\(F^1=H\) and \(F^2=0\).

\begin{proposition}\label[proposition]{prop:residual-regulator}
The assignment \(z\mapsto\cl_{\ell_0}(R_*z)\), followed by
\MathEqref{eq:single-row}, annihilates \(A_E\). It therefore
defines a map
\[
 a_{\ell_0,E}:G_E\longrightarrow H^1(E,V_{\ell_0}).
\]
The map is compatible with restriction, corestriction and
correspondences preserving the indicated projectors. Under the
common cohomological identification, it is independent of the
maximal curve-visible projector \(e\).
\end{proposition}

\begin{proof}
For a degree-zero divisor \(D\) on a pointed curve \(C\), its
continuous cycle class lies in the \(\pi_{1,C}\)-part of
\(H^2_{\cont}(C,\Q_{\ell_0}(1))\). This part has only the
Hochschild--Serre term
\(H^1(E,H^1(C_{\Fbar},\Q_{\ell_0}(1)))\).
If \(\Delta\in\CH^c(C\times Y)_\Q\) is a projected curve
correspondence, \(R\Delta\pi_{1,C}\) has zero geometric
\(H^1\) action by \Cref{prop:central-projector}. By functoriality
its map on the only associated graded term is zero. Its image in
the target of \MathEqref{eq:single-row} is consequently in
\(F^2=0\).

Curve generation proves the vanishing for algebraically trivial
cycles after a finite extension \(L/E\). It descends to \(E\)
because
\(\operatorname{cor}_{L/E}\operatorname{res}_{L/E}=[L:E]\)
on rational cohomology. Restriction is therefore injective.
Compatibility with restriction, trace and the stated
correspondences follows from the cycle map and the spectral
sequence.

For independence, let \(e'\) be another maximal curve-visible
lift. Its geometric action is the same orthogonal projector onto
\(W\). The \(\Pi\)-part of total degree \(2c=k+1\) also has
only its first associated graded term. Thus equality of the
geometric actions of \(e,e'\) gives equality of their actions on
this whole continuous-cohomology group. The cycle classes of
\((\Pi-e)z\) and \((\Pi-e')z\) coincide.
\end{proof}

\subsection{Arithmetic models and weight-minus-one cohomology}
\label{sec:fixed-model}

In the notation \(G_{E,S_E}\), all archimedean places are included
in \(S_E\); displayed conditions on \(S_E\) specify its finite
places.

\begin{proposition}\label[proposition]{prop:fixed-S}
There is a finite set \(S_0\) of finite places of \(F_0\) such
that, for every prime \(\ell_0\) and finite \(E/F_0\),
\[
 a_{\ell_0,E}(G_E)
   \subseteq H^1(G_{E,S_E},V_{\ell_0})
   \subseteq H^1(E,V_{\ell_0}),
\]
where \(S_E\) consists of places above \(S_0\) or \(\ell_0\).
The first space is finite-dimensional over \(\Q_{\ell_0}\).
\end{proposition}

\begin{proof}
Spread \(Y\) to a smooth proper scheme over
\(\Spec\cO_{F_0}[1/S_0]\). Spread the finitely many cycle
representatives and rational-equivalence witnesses defining the
relative correspondence \(R\) and the identity \(R^2=R\).
After enlarging \(S_0\) once, these data all exist on this model.
For a chosen \(\ell_0\), work over
\[
 U_{\ell_0}=\Spec\cO_{F_0}
                [1/(S_0\cup\{v:v\mid\ell_0\})].
\]
The higher direct images of the smooth proper morphism are lisse.
Since their \(R\)-parts vanish at the geometric generic point in
degrees other than \(k\), they vanish everywhere in those
degrees. The surviving part is a lisse sheaf \(\cV_{\ell_0}\).
Applying \(R\) to Leray yields
\[
 R_*H^{2c}_{\cont}(\mathcal Y,\Q_{\ell_0}(c))
          \simeq H^1(U_{\ell_0},\cV_{\ell_0}).
\]
Every generic-fibre cycle extends by closure to this regular
model. Compatibility of cycle classes with restriction proves the
factorization. Arithmetic étale cohomology finiteness gives finite
\(\Q_{\ell_0}\)-dimension. Equivalently one may apply the model
argument of \cite[Proposition~2.5 and Corollary~2.6]{Raskind}.

For a finite extension \(E/F_0\), base change this same model to
\(\cO_E[1/S_E]\). This is a finite flat base change, preserves
smoothness, and preserves the relative correspondence identities.
Taking closures of cycles over \(E\) proves the assertion
uniformly. In degree one, the étale cohomology of the base with
this lisse sheaf identifies with \(H^1(G_{E,S_E},V_{\ell_0})\).
\end{proof}

The residual representation has a further useful property: its
geometric weight is \(2c-1-2c=-1\). This strengthens the preceding
image statement to one about the whole target.

\begin{proposition}[Weight-minus-one cohomology]
\label[proposition]{prop:weight-minus-one}
Let \(S_0\) contain the bad-reduction places of \(Y/F_0\).
For every finite \(E/F_0\), every prime \(\ell_0\), and every
finite \(v\notin S_E\),
\[
 H^1(E_v,V_{\ell_0})=0.
\]
Inflation is an isomorphism
\[
 H^1(G_{E,S_E},V_{\ell_0})\xrightarrow{\sim}
 H^1(E,V_{\ell_0}).
\]
In particular, the full group on the right is finite-dimensional.
For every stable free \(\Z_{\ell_0}\)-lattice
\(T\subset V_{\ell_0}\), inflation also identifies
\(H^1(G_{E,S_E},T)\) with \(H^1(E,T)\), and these groups
are finitely generated over \(\Z_{\ell_0}\).
\end{proposition}

\begin{proof}
At \(v\notin S_E\), smooth proper base change makes inertia act
trivially on \(V_{\ell_0}\). The eigenvalues of geometric
Frobenius have complex absolute value \(q_v^{-1/2}\), by
\cite[Theorem~1.6]{DeligneWeil}. Consequently
\(V_{\ell_0}/(\mathrm{Fr}_v-1)V_{\ell_0}=0\).
The \(\ell_0\)-part of tame inertia is \(\Z_{\ell_0}(1)\),
and continuous homomorphisms to \(V_{\ell_0}\) kill wild inertia
and the other tame factors. Therefore
\[
 H^1(I_v,V_{\ell_0})^{\mathrm{Fr}_v}
   =V_{\ell_0}(-1)^{\mathrm{Fr}_v}=0,
\]
since the eigenvalues on the twist have absolute value
\(q_v^{1/2}\). Local inflation--restriction proves the first
assertion.

For a global continuous cocycle, the restriction to every such
inertia group is zero. Indeed its local class is zero, and
inertia acts trivially, so even the local coboundaries restrict to
zero on inertia. The kernel of
\(\Gal(\Fbar/E)\to G_{E,S_E}\) is the closed normal subgroup
generated by these inertia groups. It acts trivially on
\(V_{\ell_0}\). If \(z\) is the cocycle and \(i\) is in one
of these inertia groups, the cocycle identity gives
\(z(gig^{-1})=(1-gig^{-1})z(g)=0\).
The restriction of \(z\) to the kernel is thus a continuous
homomorphism vanishing on all generators, hence on the whole
kernel. It factors through \(G_{E,S_E}\). Inflation is
injective, proving the isomorphism.

For \(T\), the inertia invariants are
\(T(-1)^{\mathrm{Fr}_v}=0\), as they inject into the rational
ones. The same global cocycle argument proves integral inflation.
The local integral group may be the nonzero finite group
\(T/(\mathrm{Fr}_v-1)T\); its restriction to inertia is still
zero. Finally, in degree one the continuous-cohomology convention
agrees with these continuous cocycles: compactness gives a bounded
image in a scalar multiple of a stable lattice.

To prove finite generation, write \(G=G_{E,S_E}\). The groups
\(H^i(G,T/\ell_0^aT)\) are finite by
\cite[Corollary~I.4.15]{Milne}. The inverse system of degree-zero
groups is Mittag--Leffler, so
\[
 H^1(G,T)=\varprojlim_a H^1(G,T/\ell_0^aT)
\]
is a compact pro-\(\ell_0\) module. The exact sequence for
\(0\to T\xrightarrow{\ell_0}T\to T/\ell_0T\to0\)
embeds its reduction modulo \(\ell_0\) in the finite group
\(H^1(G,T/\ell_0T)\). The compact form of Nakayama's lemma
therefore makes \(H^1(G,T)\) finitely generated. Every continuous
cocycle into \(V_{\ell_0}\) takes values in
\(\ell_0^{-a}T\) for some \(a\), and a coboundary can likewise
be cleared of denominators. Hence
\[
 H^1(G,V_{\ell_0})=H^1(G,T)\otimes_{\Z_{\ell_0}}\Q_{\ell_0},
\]
which proves the rational finite-dimensionality.
\end{proof}

\Cref{prop:residual-regulator,prop:weight-minus-one} prove
\Cref{thm:main-regulator}. These ramification statements concern
cohomology over \(E\), rather than the ramification of
\(E/F_0\). The latter enters separately in
\MathSecref{sec:fields}.

\begin{corollary}\label[corollary]{cor:finite-rank}
If, for one prime \(\ell_0\), the scalar-extended map
\begin{equation}\label{eq:scalar-detection}
 G_E\otimes_\Q\Q_{\ell_0}
       \longrightarrow H^1(E,V_{\ell_0})
\end{equation}
is injective, then \(\dim_\Q G_E<\infty\).
\end{corollary}

\begin{proof}
Its image lies in the finite-dimensional space of
\Cref{prop:fixed-S}. Scalar extension preserves the dimension of
a rational vector space.
\end{proof}

The scalar extension in \MathEqref{eq:scalar-detection} is
essential: a finite-dimensional \(\Q_{\ell_0}\)-space has
infinite dimension over \(\Q\).

\subsection{Local integral bounds}\label{sec:local-integral}

\begin{proposition}\label[proposition]{prop:local-lattice}
For every fixed \(E,\ell_0\), the subgroup
\(a_{\ell_0,E}(\Lambda_G(E))\) is contained in a finitely
generated \(\Z_{\ell_0}\)-module in the rational cohomological
target. If \MathEqref{eq:scalar-detection} is injective, then
\(\Lambda_G(E)\otimes_\Z\Z_{\ell_0}\) is finitely generated
over \(\Z_{\ell_0}\), and
\(\Lambda_G(E)/\ell_0\Lambda_G(E)\) is finite.
\end{proposition}

\begin{proof}
Choose one integer \(M\) clearing the Chow denominators of \(R\).
For an integral representative \(Z\),
\(M\cl_{\ell_0}(RZ)=\cl_{\ell_0}((MR)Z)\) is in the image of
the integral continuous cycle class map. This image is a finitely
generated \(\Z_{\ell_0}\)-module by
\cite[Theorem~2.4]{Kahn}. Let \(M_{\ell_0}\) be its image in
rational continuous cohomology. The intersection
\[
 M_{\ell_0}\cap R_*H^{2c}_{\cont}(Y_E,\Q_{\ell_0}(c))
\]
is again finitely generated and contains
\(M\cl_{\ell_0}(RZ)\). Applying the single-row identification
to this intersection and dividing by \(M\) gives the claimed bound.

Under scalar-extended injectivity, the induced map on
\(\Lambda_G(E)\otimes\Z_{\ell_0}\) is injective. Its source is
torsion-free over \(\Z_{\ell_0}\), and its localization is
\(G_E\otimes\Q_{\ell_0}\), because \(\Lambda_G(E)\) spans
\(G_E\). Its image is a submodule of a finitely generated
\(\Z_{\ell_0}\)-module, hence is finitely generated. Reduction
modulo \(\ell_0\) proves the last assertion.
\end{proof}

\section{Integral discreteness and height gaps}\label{sec:lattices}

\subsection{Real positivity forced by finiteness}\label{sec:northcott-real}

\begin{lemma}\label[lemma]{lem:northcott-real}
Let \(\Lambda\) be a subgroup spanning a rational vector space
\(V\), and let \(q\) be a real-valued quadratic form on \(V\).
If \(q\) has Northcott on \(\Lambda\), its extension to
\(V\otimes_\Q\R\) is positive definite.
\end{lemma}

\begin{proof}
It suffices to consider a finite-dimensional rational subspace
and choose a rational basis \(v_1,\ldots,v_s\) in \(\Lambda\).
Let \(Q\) be the Gram matrix. A negative real direction would,
by rational approximation and clearing denominators, give a
nonzero \(v\in\sum_i\Z v_i\) with \(q(v)<0\). Its integer
multiples contradict Northcott. Thus \(Q\) is positive semidefinite.

If \(0\ne u\in\ker Q\), round the coordinates of \(tu\) to
integers, obtaining \(z_t\in\Z^s\) with uniformly bounded
\(z_t-tu\). Then
\[
 q\Bigl(\sum_i(z_t)_iv_i\Bigr)
       =(z_t-tu)^tQ(z_t-tu)
\]
is uniformly bounded. The rounded vectors are unbounded as
\(t\to\infty\), hence take infinitely many distinct values.
This again contradicts Northcott. Therefore every such finite
Gram matrix is positive definite, which proves the assertion.
\end{proof}

In particular, residual integral Northcott already contains the
real-positivity assertion for \(q_G\). The converse requires an
arithmetic condition on the integral image.

\subsection{Finite-rank lattices}\label{sec:finite-rank-lattices}

\begin{theorem}\label[theorem]{thm:lattice-equivalence}
Let \(\Lambda\) be torsion-free of finite rank,
\(V=\Lambda\otimes\Q\), and suppose that \(q\) is real positive
definite on \(V\otimes\R\). The following are equivalent:
\begin{enumerate}[label=(\roman*)]
\item \(q\) has Northcott on \(\Lambda\).
\item There is \(\varepsilon>0\) such that
\(q(\lambda)\ge\varepsilon\) for every nonzero \(\lambda\in\Lambda\).
\item \(\Lambda\) is finitely generated.
\item For a rational basis \(v_1,\ldots,v_s\) of \(V\), there is
an integer \(N\ge1\) such that
\(\Lambda\subset N^{-1}\sum_i\Z v_i\).
\item The diagonal image of \(\Lambda\) in \(V(\A_f)\) is
relatively compact.
\end{enumerate}
The zero group satisfies these statements with the gap condition
interpreted vacuously.
\end{theorem}

\begin{proof}
Finiteness in a unit height ball gives a positive minimum, or the
bound \(1\) if that ball contains no nonzero point. This proves
(i) implies (ii). The gap makes \(\Lambda\) uniformly discrete
for the Euclidean norm \(\sqrt q\). Choose a rational basis in
\(\Lambda\) and reduce elements of \(\Lambda\) modulo its
\(\Z\)-span into a bounded fundamental parallelepiped. Uniform
discreteness makes the set of representatives finite. Hence this
span has finite index in \(\Lambda\), proving (iii) and (iv).
Conversely, a bounded subset of a Euclidean lattice is finite,
so (iv) implies (i). Finite generation gives (iv) for any rational
basis by clearing the finitely many coefficients.

Every compact subset of \(V(\A_f)\) is contained in
\(N^{-1}\Lambda_0\otimes\widehat\Z\) for some integer \(N\)
and a fixed reference lattice \(\Lambda_0\). This follows by a
finite subcover from the open compact sets of that form. Conversely
each such set is compact. Intersecting with \(V(\Q)\) proves the
equivalence with (v).
\end{proof}

If \(\Lambda\subset N^{-1}\Lambda_0\), and the least eigenvalue
of the Gram matrix in a basis of \(\Lambda_0\) is \(\mu>0\),
one obtains the explicit estimate
\begin{equation}\label{eq:lattice-count}
 \#\{\lambda\in\Lambda:q(\lambda)\le B\}
 \le\left(1+2\left\lfloor N\sqrt{B/\mu}\right\rfloor\right)^s
 \qquad(B\ge0).
\end{equation}

\subsection{An all-prime saturation criterion}\label{sec:all-prime}

\begin{theorem}\label[theorem]{thm:all-prime}
Let \(V\) be a finite-dimensional rational vector space,
\(\Lambda_0\subset V\) a full lattice, and \(\Lambda\subset V\)
a subgroup. For every prime \(p\), let
\[
 D_p:V\otimes\Q_p\hookrightarrow W_p
\]
be an injective linear map to a finite-dimensional space, and let
\(M_p\subset W_p\) be a full \(\Z_p\)-lattice. Assume
\[
 D_p(\Lambda)\subseteq M_p\quad\text{for every }p,
 \qquad D_p^{-1}(M_p)=\Lambda_0\otimes\Z_p
               \quad\text{for almost every }p.
\]
Then \(\Lambda\subset N^{-1}\Lambda_0\) for one integer \(N\).
Every real positive definite quadratic form has Northcott on
\(\Lambda\).
\end{theorem}

\begin{proof}
Injectivity makes \(L_p=D_p^{-1}(M_p)\) a \(\Z_p\)-lattice:
it is open and compact in \(V\otimes\Q_p\). For each exceptional
prime choose \(b_p\ge0\) with
\(L_p\subset p^{-b_p}\Lambda_0\otimes\Z_p\), and put
\(N=\prod p^{b_p}\). The coordinates of \(N\lambda\) lie in
\(\Z_p\) for every \(p\). Since
\(\Q\cap\prod_p\Z_p=\Z\), they are integers. Apply
\Cref{thm:lattice-equivalence}.
\end{proof}

For the residual regulator, a common multiple of the projector
denominators can first be absorbed into \(D_p\). Once scalar
injectivity is available, \Cref{prop:local-lattice} supplies local
integral modules. The almost-everywhere saturation condition in
\Cref{thm:all-prime} is the additional comparison with one common
rational lattice.

\begin{example}\label[example]{ex:denominators}
Let
\[
 \Lambda=\sum_{p\ {
m prime}}\Z\,\frac1p\subset\Q,
 \qquad q(x)=x^2.
\]
Its elements have squarefree denominators. For every prime \(p\),
\(\Lambda\otimes\Z_p=p^{-1}\Z_p\), which is finite free of
rank one. Nevertheless \(q(1/p)=p^{-2}\to0\), so Northcott
fails. The integral maps \(D_p(x)=px\) have inverse lattices
larger than \(\Z_p\) at infinitely many primes. Thus local
finite generation at every prime does not provide the common
denominator required by \Cref{thm:lattice-equivalence}.
\end{example}

Already \(\Z[1/p]\), with \(q(x)=x^2\), defeats any finite
collection of integral detectors at primes different from \(p\).
Another distinction is real definiteness: the form
\((a-\sqrt2b)^2\) on \(\Z^2\) is positive on nonzero rational
vectors, but rounding along its real kernel produces infinitely
many bounded-height lattice vectors.

\subsection{A quantitative descent criterion}\label{sec:weak-mw}

\begin{theorem}\label[theorem]{thm:weak-mw-gap}
Let \(\Lambda\) be torsion-free and let \(\sqrt q\) be the norm
of a real positive quadratic form on \(\Lambda\otimes\R\).
Suppose, for some integer \(m\ge2\),
\[
 s=|\Lambda/m\Lambda|<\infty,\qquad
 q(\lambda)\ge\varepsilon>0\quad(0\ne\lambda\in\Lambda).
\]
Then \(\Lambda\) is finitely generated and \(q\) has Northcott.
More precisely, if \(B\ge0\), \(j\in\Z_{\ge0}\), and
\(m^{2j}\varepsilon>4B\), then
\begin{equation}\label{eq:kummer-count}
 \#\{\lambda\in\Lambda:q(\lambda)\le B\}\le s^j.
\end{equation}
\end{theorem}

\begin{proof}
If \(x,y\) are in the height-\(B\) ball and
\(x-y=m^jz\), then
\(\|z\|\le2\sqrt B/m^j<\sqrt\varepsilon\), so \(z=0\).
Reduction modulo \(m^j\Lambda\) is injective on that ball.
Torsion-freeness identifies the successive quotients
\(m^i\Lambda/m^{i+1}\Lambda\) with \(\Lambda/m\Lambda\),
proving \MathEqref{eq:kummer-count}.

Choose a finite set \(T\) of representatives modulo
\(m\Lambda\), and let \(C=\max_{t\in T}\|t\|\).
Writing \(x=my+t\) gives
\(\|y\|\le(\|x\|+C)/m\). Repetition brings \(y\) into
any fixed ball of radius greater than \(C/(m-1)\). This ball
is finite by the first part; together with \(T\), it generates
\(\Lambda\). This is the height-descent argument underlying
weak Mordell--Weil; compare \cite[Section~10.3]{ConradFeng}.
\end{proof}

Northcott by itself need not give finite rational rank. For example,
\(\Lambda=\bigoplus_{i\ge1}\Z e_i\) with
\(q(\sum a_ie_i)=\sum2^ia_i^2\) has finite height balls and
infinite rank. Rational detection, integral discreteness and real
positivity therefore retain distinct roles in the general program.

The examples in this section are examples of abstract groups and
quadratic forms. Their purpose is to identify the exact arithmetic
hypotheses in a proof of geometric Northcott.

\section{Fields of classes and bounded-degree finiteness}\label{sec:fields}

\subsection{Descent by trace}\label{sec:class-fields}

Let \(\mathcal G\) be a rational geometric Chow quotient with its
discrete \(\Gal(\Fbar/F_0)\)-action, and let \(\Lambda(E)\subset\mathcal G\)
be compatible integral images. We assume that restriction and trace
preserve these images in their usual directions. Fixed rational Chow
projectors and quotient by geometric algebraic equivalence have this
property. Every element considered is defined over a finite extension,
so its stabilizer is open. Define its minimal class field by
\[
 F_0(\xi)=\Fbar^{\operatorname{Stab}_{\Gal(\Fbar/F_0)}(\xi)}.
\]

\begin{lemma}\label[lemma]{lem:class-trace}
If \(\xi\in\Lambda(E)\) and \(K=F_0(\xi)\), then
\[
 K\subseteq E,\qquad [E:K]\xi\in\Lambda(K).
\]
In particular, if \([E:F_0]\le D\), the required multiplier is
at most \(D\).
\end{lemma}

\begin{proof}
The inclusion \(\Gal(\Fbar/E)\subseteq\operatorname{Stab}(\xi)\) gives
\(K\subseteq E\). Take an integral cycle over \(E\) representing
\(\xi\) in the indicated quotient and push it forward to \(K\).
Over \(\Fbar\), the resulting class is the sum over the
\(K\)-embeddings of \(E\). Every summand equals \(\xi\), so
the sum is \([E:K]\xi\). Fixed projectors commute with trace,
and algebraic equivalence is preserved by it.
\end{proof}

The lemma concerns the field of a class in a rational quotient.
It supplies the precise denominator even when no chosen integral
cycle representative descends to that field.

\subsection{The ramification criterion}\label{sec:ramification-proof}

Let \(q\) be a Galois-invariant quadratic form on \(\mathcal G\).
For \(B\ge0\) and \(D\ge1\), consider
\[
 \mathcal B(B,D)=
 \{\xi\in\bigcup_{[E:F_0]\le D}\Lambda(E):q(\xi)\le B\}.
\]

\begin{theorem}\label[theorem]{thm:ramification}
Assume \(q\) has fixed-field Northcott on \(\Lambda(K)\) for
every finite \(K/F_0\). Then bounded-degree Northcott is
equivalent to the following condition: for every \(B,D\), there
is a finite set \(S(B,D)\) of finite places of \(F_0\) such that
each field \(F_0(\xi)/F_0\), with \(\xi\in\mathcal B(B,D)\),
is unramified outside \(S(B,D)\).
\end{theorem}

\begin{proof}
Suppose the ramification condition holds. By the finite-ramification
form of Hermite--Minkowski, there are only finitely many extensions
\(K/F_0\) in \(\Fbar\) of degree at most \(D\) unramified
 outside \(S(B,D)\); see \cite[Theorem~10.2.2]{Kedlaya}.
This form also follows by bounding the local discriminant exponents
in fixed degree and applying discriminant finiteness.

For \(\xi\in\mathcal B(B,D)\), \Cref{lem:class-trace} gives
\(m\xi\in\Lambda(K)\), where \(K=F_0(\xi)\) is among these
finitely many fields and \(1\le m\le D\). Moreover
\[
 q(m\xi)=m^2q(\xi)\le D^2B.
\]
Fixed-field Northcott leaves finitely many possibilities for
\(m\xi\). Multiplication by \(m\) is injective in the rational
group, and only finitely many \(m\) occur. Thus the ball is finite.
Conversely, if the ball is finite, the union of the ramification
sets of its finitely many class fields is finite.
\end{proof}

An explicit discriminant estimate would give a quantitative route
to the hypothesis. For example, with \(\mathfrak d_{K/F_0}\)
the relative discriminant ideal, an inequality
\begin{equation}\label{eq:discriminant-target}
 \log N_{F_0/\Q}\mathfrak d_{F_0(\xi)/F_0}
       \le A_D q(\xi)+B_D
\end{equation}
on integral classes of degree at most \(D\) would suffice.
The qualitative ramification condition only asks for some finite
set depending on \(B,D\); it is weaker than
\MathEqref{eq:discriminant-target}.

\subsection{Why uniform fixed-field estimates do not suffice}
\label{sec:quadratic-example}

\begin{example}\label[example]{ex:quadratic-fields}
Choose infinitely many distinct quadratic extensions \(K_i/F_0\)
with characters \(\chi_i\). Put
\[
 \mathcal G=\bigoplus_i\Q e_i,\qquad
 \sigma e_i=\chi_i(\sigma)e_i,\qquad
 \Lambda(E)=\bigoplus_{K_i\subseteq E}\Z e_i,
 \qquad q\left(\sum a_ie_i\right)=\sum a_i^2.
\]
The lattices are compatible with restriction and trace. Every
finite extension \(E\) has at most \([E:F_0]-1\) distinct
quadratic subfields: their compositum is multiquadratic, and the
number of its quadratic subfields is one less than its degree.
Consequently
\[
 \rk\Lambda(E)\le[E:F_0]-1,\qquad
 q(\lambda)\ge1\ (\lambda\ne0),\qquad
 |\Lambda(E)/m\Lambda(E)|\le m^{[E:F_0]-1}.
\]
All fixed-field height balls are finite. Nevertheless the classes
\(e_i\in\Lambda(K_i)\) all have degree two and height one.
Thus bounded-degree Northcott fails despite uniform
degree-dependent bounds on ranks, gaps and Kummer quotients.
\end{example}

The fixed-\(S\) cohomological statement of
\Cref{prop:fixed-S} holds over a chosen field of definition.
\Cref{thm:ramification} instead bounds the fields themselves.
This distinction is essential when using residual arithmetic
regulators to study the conjectural Northcott property.

\section{Normalized determinants on Jacobian families}\label{sec:metrics}

\subsection{The Poincaré metric and normalization}\label{sec:poincare-metric}

Let \(C\) be a smooth projective connected complex curve with
basepoint \(o\), and let \(J=\Pic^0(C)\). Let \(\cP\) be the
Poincaré line bundle on \(C\times J\), rigidified on both axes.
Equip it with the standard smooth metric for which these
rigidifications have norm one. Its Chern form
\(\Omega=c_1(\cP)\) is purely mixed: it has real degree one in
\(C\) and degree one in \(J\).

One way to see this is to pull back the canonical Poincaré bundle
on \(J\times J^\vee\) by the Abel map and the principal
polarization, with the convention realizing \(\cP\). For a
translation-invariant theta curvature \(\theta\), the corresponding
Poincaré curvature is the mixed form obtained from
\(m^*\theta-p_1^*\theta-p_2^*\theta\), up to the chosen duality
sign. The pure terms cancel. Pullback by the Abel map preserves
the mixed-degree assertion.

For a fixed smooth projective complex variety \(Z\), use the
cochain determinant convention
\[
 \lambda_Z(K)=\det R\Gamma(Z,K)
       =\bigotimes_i\det H^i(Z,K)^{(-1)^i}.
\]
A perfect complex is given a bounded vector-bundle representative
with smooth hermitian metrics; its determinant metric is the
alternating tensor product of the Quillen metrics of the terms.
Choose fixed Kähler metrics on the varieties involved.

\subsection{The normalized projection-formula isometry}
\label{sec:quillen-proof}

Let \(X\) be a smooth projective connected complex variety,
\(M=C\times X\), and \(q:M\to C\). For a hermitian perfect
complex \(\mathcal F\) on \(M\), put \(K=Rq_*\mathcal F\)
and choose any hermitian vector-bundle representative of \(K\).
Define line bundles on \(J\) by
\begin{align*}
 \cD_M(\mathcal F)
  &=\det Rp_{J*}(\cP_M\otimes\mathcal F_J)
     \otimes\bigl(\lambda_M(\mathcal F)\otimes\cO_J\bigr)^{-1},\\
 \cD_C(K)
  &=\det Rp_{J*}(\cP\otimes K_J)
     \otimes\bigl(\lambda_C(K)\otimes\cO_J\bigr)^{-1}.
\end{align*}
The subscripts indicate pullbacks to the appropriate products with
\(J\). Each constant denominator has the Quillen norm of the same
complex and metrics that occur in the numerator at \(0\).
Thus each line has a canonical origin rigidification of norm one.

\begin{theorem}\label[theorem]{thm:quillen-isometry}
Use product Kähler metrics on \(C\times X\). The derived
projection-formula isomorphism gives an exact isometry of
rigidified line bundles
\[
 \cD_M(\mathcal F)\xrightarrow{\sim}\cD_C(Rq_*\mathcal F).
\]
Both Chern forms equal
\begin{equation}\label{eq:quillen-curvature}
 \frac{\rk K}{2}\int_C\Omega^2.
\end{equation}
Under the canonical rigidified identifications, these normalized
metrics are independent of the hermitian representatives and the
fixed product Kähler metrics used to construct them.
\end{theorem}

\begin{proof}
Derived projection formula and base change identify the numerator
complexes and their restrictions at \(0\). Applying determinants
to these two identifications gives the asserted algebraic
isomorphism, respecting the origin rigidifications.

The Bismut--Gillet--Soulé curvature formula
\cite[Theorem~0.1]{BGS}, in the cochain determinant convention,
gives the Chern form of the first line as the base-degree-two
component of
\[
 \int_{C\times X}\Td(T_C)\Td(T_X)\ch(\mathcal F)
                          (e^\Omega-1).
\]
Every factor except \(\Omega\) is pulled back from \(C\times X\).
Since \(\Omega\) is purely mixed, only \(\Omega^2/2\) can
contribute degree two on \(J\). It already has degree two on
\(C\), so the remaining factor contributes its degree-zero-in-\(C\)
part. Integration over \(X\) gives
\[
 \int_X\Td(T_X)\ch(\mathcal F|_{\{x\}\times X})
       =\chi(X,\mathcal F|_{\{x\}\times X})=\rk K
 \qquad(x\in C).
\]
The final equality follows by derived base change and the constant
Euler characteristic of the perfect pushforward. This proves
\MathEqref{eq:quillen-curvature} for \(\cD_M\). The curve
calculation gives the same form: all positive-degree forms from
\(K\) or \(T_C\) have too much \(C\)-degree to contribute.

The determinant lines with their Quillen metrics are smooth on
the whole \(J\), including cohomology-jump loci. Under the
algebraic isomorphism, the logarithm of the ratio of their metrics
is therefore a smooth pluriharmonic function on the compact
connected manifold \(J\). It is constant, and the norm-one
origin rigidifications make that constant zero. Changing the
hermitian representatives or the fixed product metrics preserves both the
curvature computation and this normalization, proving the last
assertion.
\end{proof}

\subsection{Virtual rank zero and Deligne pairings}
\label{sec:deligne-proof}

\begin{theorem}\label[theorem]{thm:deligne-isometry}
Let \(K\) be a perfect complex of virtual rank zero on \(C\),
and put \(L=\det K\). There is a unique rigidified isomorphism
\[
 \cD_C(K)\xrightarrow{\sim}
       \langle\cP,p_C^*L\rangle_{C\times J/J}.
\]
It is an isometry for the normalized Quillen metric and the
Deligne metric. Both metrics are smooth and flat. The right side
can be formed using any smooth hermitian metric on \(L\).
\end{theorem}

\begin{proof}
On a nonsingular curve, rank and determinant identify
\(K_0(C)\) with \(\Z\oplus\Pic(C)\). Thus
\([K]=[L]-[\cO_C]\). At the level of line-bundle isomorphism
classes, the normalized determinant becomes
\[
 \lambda_C(\cP L)\lambda_C(\cP)^{-1}
       \lambda_C(L)^{-1}\lambda_C(\cO_C),
\]
the determinant expression for the Deligne pairing. Here the first
two factors are relative determinants over \(J\). Normalize an
isomorphism at \(0\). It becomes unique because an invertible
holomorphic function on the connected projective \(J\) is
constant.

The Quillen curvature is zero by
\MathEqref{eq:quillen-curvature}. The Deligne curvature is
\(\int_C\Omega\,c_1(L)=0\), since the integrand has degree
three on \(C\). At \(0\), \(\cP\) is the trivially metrized
line with unit rigidification; the Deligne unit has norm one.
The two flat metrics therefore agree by the same compactness
argument as above. The determinant and Deligne curvature formulas
are also discussed in \cite[Section~2]{BiswasSchumacher}.
\end{proof}

\subsection{Curve correspondences in arbitrary codimension}
\label{sec:correspondence-metric}

Let \(\dim X=n\), take
\(\Gamma\in\CH^r(C\times X)_\Q\), and let
\(W\in V^{n+1-r}(X)\). Choose rational \(K_0\)-lifts
\(E_\Gamma,E_W\) whose Chern characters start with \(\Gamma,W\)
in their indicated codimensions. Virtual sums of structure sheaves
give such lifts on these smooth varieties. Set
\[
 \mathcal F=E_\Gamma\otimes^{\mathbf L}p_X^*E_W,
 \qquad K=Rq_*\mathcal F.
\]
The product Chern character starts in codimension \(n+1\).
Filtered Grothendieck--Riemann--Roch for \(q\) gives
\begin{equation}\label{eq:grr-leading}
 \rk K=0,\qquad
 c_1(\det K)=q_*(\Gamma\cdot p_X^*W)={}^t\Gamma_*W.
\end{equation}
Positive-degree Todd terms do not contribute to these two
components. The divisor class on the right has degree zero because
\(W\) is homologically trivial.

\begin{corollary}\label[corollary]{cor:cycle-determinant}
There is a rigidified isometry of rational metrized line bundles
on \(J\),
\[
 \cD_{C\times X}(E_\Gamma\otimes^{\mathbf L}p_X^*E_W)
   \xrightarrow{\sim}
 \langle\cP,\cO_C({}^t\Gamma_*W)\rangle.
\]
The metric is the normalized flat Poincaré metric attached to the
degree-zero divisor class \({}^t\Gamma_*W\).
\end{corollary}

\begin{proof}
Combine \Cref{thm:quillen-isometry,thm:deligne-isometry} with
\MathEqref{eq:grr-leading}. For rational classes, take a common
tensor power throughout and then the positive real root of the
norm. Changing the chosen cycle lifts by higher filtration changes
their product only in codimension at least \(n+2\), which is
zero on \(C\times X\); the rational Chern character then gives
the same normalized line.
\end{proof}

The relation to the classical archimedean height uses the
determinant--cycle--Hodge biextension comparison of
\cite[Propositions~4.3, 4.29, 4.31 and Corollary~4.32]{Gorchinskiy}.
Its disjoint-cycle trivialization gives the local height
\cite[Remark~4.7]{Gorchinskiy}. The relative comparison over a
smooth base is given in \cite[Theorem~10.9]{Nakayama}; that theorem
includes an inversion of the central \(\mathbf G_m\) under its
conventions. Fixing the cochain determinant convention above and
the corresponding Poincaré convention fixes this sign.

An incidence locus can be a zero or pole locus of the chosen
meromorphic frame without being a singularity of the line metric.
If \(s=fe\), where \(e\) is a nonvanishing holomorphic frame, then
\[
 -\log\|s\|=-\log|f|-\log\|e\|.
\]
The second term is smooth across incidence in the Jacobian family
just constructed. This global line comparison resolves the
normalization and continuity issue for that determinant model.

\subsection{The source-curve good-reduction comparison}
\label{sec:good-reduction-metric}

Let \(\mathcal O\) be a complete discrete valuation ring with
fraction field \(K_v\), equipped with a fixed non-Archimedean
absolute value inducing the given discrete valuation.
Let \(\mathcal C/\mathcal O\) be a smooth proper geometrically
connected curve with a section, let \(\mathcal J/\mathcal O\)
be its Jacobian, and let \(\mathcal P\) be the axis-rigidified
Poincaré bundle. For a perfect complex \(K\) on \(\mathcal C\),
form the normalized determinant line \(\cD_{\mathcal C}(K)\)
on \(\mathcal J\) by the same algebraic construction.

\begin{proposition}\label[proposition]{prop:good-reduction}
Suppose the generic virtual rank of \(K\) is zero and
\(\deg\det K_{K_v}=0\). The model metric of
\(\cD_{\mathcal C}(K)\) on the generic fibre is the canonical
rigidified Poincaré metric associated with \(\det K_{K_v}\).
\end{proposition}

\begin{proof}
The degree of \(\det K\) is zero on every fibre. After
rigidification along the section, it defines a point of the
relative \(\Pic^0\), and hence a rigidified Poincaré extension
on \(\mathcal J\). The rank-zero determinant identity on the
generic curve identifies its generic line with
\(\cD_{\mathcal C}(K)\), respecting the origin rigidification.

Regard this isomorphism as a rational section of the quotient of
the two model lines. Its divisor is vertical. The abelian scheme
\(\mathcal J\) is regular and its special fibre is geometrically
integral, so this divisor is \(a[\mathcal J_s]\) for an integer
\(a\). Pullback to the zero section has valuation \(a\), but
the two integral unit rigidifications make that valuation zero.
Thus \(a=0\). Normality extends the isomorphism and its inverse
over the whole model. The model metrics coincide. The Poincaré
multiplication laws extend over the abelian scheme and are
isometries, identifying this model metric with the canonical one.
\end{proof}

If an ambient model \(\mathcal X\) and a model-perfect complex
on \(\mathcal C\times\mathcal X\) have a perfect derived
pushforward to \(\mathcal C\), the determinant projection
formula reduces their model comparison to this proposition.
At a bad-reduction source curve the vertical correction involves
the geometry of a semistable or Néron model; this is a separate
comparison problem.

\section{An equivalent conjectural package}\label{sec:program}

\subsection{The three central assertions}\label{sec:three-assertions}

Fix a product \(X=\prod_{i=1}^nC_i\), its points \(o_i\), and
\(\ell=\sum_i a_iH_i\), with \(a_i>0\). Pass to one field
\(F_0\) defining the maximal curve-visible projectors for all
odd subproducts. All height assertions use
\Cref{def:heights}.

\begin{conjecture}[Off-central vanishing]
\label[conjecture]{conj:offcentral}
For every subproduct \(Y_S\), every finite \(E/F_0\), and
every \(c\) with \(|S|\ne2c-1\),
\[
 U_S^c(E)=0.
\]
\end{conjecture}

\begin{conjecture}[Central algebraic Abel--Jacobi injectivity]
\label[conjecture]{conj:AJ}
For every odd subproduct \(|S|=2c-1\) and every finite \(E/F_0\),
the rationalized Abel--Jacobi map is injective on \(A_{S,E}\).
\end{conjecture}

\begin{conjecture}[Residual integral Northcott]
\label[conjecture]{conj:residual-northcott}
For every odd subproduct \(|S|=2c-1\), the quadratic form
\(q_{G,S}=(-1)^ch_{G,S}\) has bounded-degree Northcott on
the projected integral image \(\Lambda_{G,S}\).
\end{conjecture}

The real-positivity assertion in the central program is a
consequence of \Cref{conj:residual-northcott}, by
\Cref{lem:northcott-real}. It can still be studied first as the
weaker conjecture that \(q_{G,S}\) is real positive definite.
\Cref{ex:denominators} shows the additional integral issue in the
converse direction.

\begin{theorem}[Equivalence of the product and residual programs]
\label[theorem]{thm:program-equivalence}
For the fixed pointed product and the displayed polarization, the
following packages are equivalent.
\begin{enumerate}[label=(\roman*)]
\item Arithmetic hard Lefschetz holds in every degree and over
every finite extension of \(F_0\), and
\(q_{\ell,r}\) has bounded-degree Northcott on the primitive
integral image in every degree
\(1\le r\le\lfloor(n+1)/2\rfloor\).
\item \Cref{conj:offcentral,conj:AJ,conj:residual-northcott} hold
for all the subproducts of \(X\).
\end{enumerate}
Either package implies real primitive positivity and
nondegeneracy of the complementary height pairings over every
fixed number field.
\end{theorem}

\begin{proof}
From (ii), \Cref{thm:main-lefschetz} gives arithmetic hard
Lefschetz. \Cref{thm:central-northcott} gives central integral
Northcott, and \Cref{lem:northcott-real} gives real central
positivity. \Cref{cor:positivity-propagation} gives all the
primitive Northcott and positivity assertions, as well as
nondegeneracy.

Conversely, hard Lefschetz in (i) gives off-central vanishing by
\Cref{thm:main-lefschetz}. For an odd subproduct \(|S|=2c-1\),
inflate \(u\in U_S^c(E)\) by fundamental classes on all omitted
factors. In the notation of \MathSecref{sec:chow-blocks}, this is
the integral injection \(u\mapsto j_\varnothing u\) in degree
\(r=c\). With \(m=n-|S|\), it is primitive and satisfies
\[
 q_{\ell,c}(j_\varnothing u)
       =m!\prod_{i\notin S}a_i\,q_{S,c}(u).
\]
Thus primitive Northcott on \(X\) gives central integral
Northcott. \Cref{thm:central-northcott} supplies
\Cref{conj:AJ,conj:residual-northcott}. The final assertion now
follows from the first direction of the proof.
\end{proof}

This is a product-of-curves formulation of the arithmetic
Lefschetz and Hodge--Riemann program, with an integral finite-height
condition added. The underlying arithmetic standard conjectures
are those in \cite[Conjecture~2.3.1]{ZhangStandard}.

\subsection{Arithmetic detection and saturation}
\label{sec:detection-conjectures}

Fix one central residual group \(G_E\).

\begin{conjecture}[Residual rational detection]
\label[conjecture]{conj:detection}
For every finite \(E/F_0\) and every prime \(p\), the map
\[
 G_E\otimes_\Q\Q_p\longrightarrow H^1(E,V_p)
\]
of \Cref{prop:residual-regulator} is injective.
\end{conjecture}

This is a residual form of rational arithmetic cycle-class
injectivity, in the framework of \cite{Raskind,Kahn}. Together
with \Cref{thm:lattice-equivalence}, it reduces fixed-field
Northcott under real positivity to the following integral question.

\begin{conjecture}[Residual integral discreteness]
\label[conjecture]{conj:lattice}
For every finite \(E/F_0\), \(G_E\) is finite-dimensional and
there are a rational basis \(v_1,\ldots,v_s\) and an integer
\(N_E\ge1\) with
\[
 \Lambda_G(E)\subseteq N_E^{-1}\sum_i\Z v_i.
\]
\end{conjecture}

One route is to prove the common-lattice, almost-everywhere
saturation hypothesis of \Cref{thm:all-prime} for the arithmetic
regulators. A second route is a genuine positive height gap,
together with the finite Kummer quotients of
\Cref{thm:weak-mw-gap}. The local integral images of
\Cref{prop:local-lattice} are already controlled; their comparison
as the coefficient prime varies is the remaining step.

Integral detection must refer to the stated image in a rational
quotient. The unrestricted integral continuous cycle class map
has torsion counterexamples over number fields in codimension
three \cite[Theorems~1.3--1.4]{ScaviaSuzuki}. Those groups and
their torsion kernels are different from \(\Lambda_G(E)\).

\subsection{Varying fields and local metrics}\label{sec:next-targets}

\begin{conjecture}[Residual class-field ramification]
\label[conjecture]{conj:ramification}
For each \(B\ge0,D\ge1\), the minimal class fields
\(F_0(\xi)\) of \(\xi\in\Lambda_G(E)\), with
\([E:F_0]\le D\) and \(q_G(\xi)\le B\), are unramified
outside one finite set of places depending on \(Y,\Pi,B,D\).
\end{conjecture}

With fixed-field Northcott, this is exactly the remaining
bounded-degree condition by \Cref{thm:ramification}. A quantitative
version asks whether constants \(A_D,B_D\) in
\MathEqref{eq:discriminant-target} exist. The quantitative
inequality is recorded as a research question, rather than being
identified with the qualitative conjecture.

For the metric problem, the next comparison is local and precise.
Starting with a fixed semistable source curve, construct the
admissible correction to the model determinant and compare its
rigidified metric with the classical Poincaré metric over the
original valued field. \Cref{prop:good-reduction} is the smooth
model case; the complex-place comparison is
\Cref{thm:quillen-isometry,cor:cycle-determinant}. Applying these
comparisons to an actual-form non-Archimedean Green construction
requires its trace, covariance and normalization statements on
their full test spaces.

Boundary estimates in a varying family also require their own
geometric input. The smooth relative biextension comparison in
\cite[Theorem~10.9]{Nakayama} is unconditional. In that work,
the singular-fibre extension in Theorem~11.2 uses a modified cycle
satisfying a specified cohomological condition, and Lemma~11.1
obtains that modification from the Hodge conjecture. These are
the hypotheses to verify when importing boundary asymptotics
into the family-height program.

\begin{thebibliography}{99}
\addcontentsline{toc}{section}{References}

\bibitem{BGS}
J.-M. Bismut, H. Gillet and C. Soulé,
\emph{Analytic torsion and holomorphic determinant bundles I:
Bott--Chern forms and analytic torsion},
Communications in Mathematical Physics \textbf{115} (1988), 49--78.
\href{https://www.imo.universite-paris-saclay.fr/~jean-michel.bismut/Bismut/1988.pdf}
{Author-hosted PDF}.

\bibitem{BiswasSchumacher}
I. Biswas and G. Schumacher,
\emph{Deligne pairing and Quillen metric},
\href{https://arxiv.org/abs/1501.02539}{arXiv:1501.02539}.

\bibitem{ConradFeng}
B. Conrad, notes by T. Feng,
\emph{Abelian varieties}, Stanford lecture notes, Spring 2015.
\href{https://math.stanford.edu/~conrad/249CS15Page/handouts/abvarnotes.pdf}
{Author-hosted PDF}.

\bibitem{DeligneWeil}
P. Deligne, \emph{La conjecture de Weil I},
Publications Mathématiques de l'IHÉS \textbf{43} (1974), 273--307.
\href{https://www.numdam.org/article/PMIHES_1974__43__273_0.pdf}
{Original article}.

\bibitem{Gorchinskiy}
S. Gorchinskiy, \emph{Notes on the biextension of Chow groups},
\href{https://arxiv.org/abs/0802.1437v2}{arXiv:0802.1437v2}.

\bibitem{Kahn}
B. Kahn, \emph{On the injectivity and non-injectivity of the
\(\ell\)-adic cycle class maps}, manuscript dated 10 March 2024.
\href{https://webusers.imj-prg.fr/~bruno.kahn/preprints/ladic-cycle-class-rev3.pdf}
{Author-hosted PDF}.

\bibitem{Kedlaya}
K. S. Kedlaya, \emph{Weil cohomology in practice},
Chapter 10, Comparing Galois representations: the Faltings--Serre
method, lecture of 30 October 2019.
\href{https://kskedlaya.org/weil-cohom/chapter-10.html}
{Author-hosted notes}.

\bibitem{LagaShnidman}
J. Laga and A. Shnidman, \emph{Vanishing criteria for Ceresa cycles},
Compositio Mathematica, published online 29 December 2025.
\href{https://www.cambridge.org/core/journals/compositio-mathematica/article/vanishing-criteria-for-ceresa-cycles/051A1B2257FB703EA13086A238A279B8}
{Publisher's full text}.

\bibitem{Nakayama}
Y. Nakayama, \emph{Asymptotics of local height pairing},
\href{https://arxiv.org/abs/2512.22551v1}{arXiv:2512.22551v1},
27 December 2025.

\bibitem{Milne}
J. S. Milne, \emph{Arithmetic duality theorems}, second edition,
BookSurge, 2006.
\href{https://jmilne.org/math/Books/ADTnot.pdf}{Author-hosted book}.

\bibitem{Raskind}
W. Raskind, \emph{Higher \(\ell\)-adic Abel--Jacobi mappings and
filtrations on Chow groups},
\href{https://arxiv.org/abs/alg-geom/9412001}{arXiv:alg-geom/9412001}.

\bibitem{ScaviaSuzuki}
F. Scavia and F. Suzuki,
\emph{Non-injectivity of the cycle class map in continuous
\(\ell\)-adic cohomology},
\href{https://arxiv.org/abs/2208.06713}{arXiv:2208.06713}.

\bibitem{ZhangStandard}
S.-W. Zhang, \emph{Standard conjectures and height pairings},
\href{https://arxiv.org/abs/2009.07089v3}{arXiv:2009.07089v3}.
The numbering cited here refers to the version of 29 April 2022.

\end{thebibliography}

\end{document}
